\section*{Bab \ref{ch:b2:field-lab-methods-2} --- Metode Lapangan dan Laboratorium II}
% numbers from the Book 4 second-writer scratch calculator calc/ch31.py

\begin{solution}{exo:b2:field-lab-methods-2:1}
(a) (i) \omterm{def:b2:field-lab-methods-2:textures}{Tekstur ofitik}; (ii) \omterm{def:b2:field-lab-methods-2:textures}{tekstur granoblastik}; (iii) \omterm{def:b2:field-lab-methods-2:porphyroblast}{poikiloblas}, yaitu \omterm{def:b2:field-lab-methods-2:porphyroblast}{porfiroblas} yang penuh
inklusi; (iv) \omterm{def:b2:field-lab-methods-2:porphyroblast}{korona}, yaitu tepian reaksi di sekeliling garnet yang sedang terurai. (b) Menurut
\eqref{eq:b2:field-lab-methods-2:lami}, $\gamma_1/\gamma_2 = \sin\ang{104}/\sin\ang{128} =
0.9703/0.7880 = 1.23$: batas itu menarik \qty{23}{\percent} lebih kuat daripada dua batas lainnya.
\end{solution}

\begin{solution}{exo:b2:field-lab-methods-2:2}
(a) $Q + A + P = 80$: Q $= 100 \times 25/80 = 31.3$, A $= 18.8$, P $= 50.0$. (b) P/(A+P) $=
40/55 = \qty{72.7}{\percent}$, antara 65 dan 90 dengan Q antara 20 dan 60: granodiorit.
(c) $M' = 12 + 8 = \qty{20}{\percent}$, di bawah 35: leukokratik.
\end{solution}

\begin{solution}{exo:b2:field-lab-methods-2:3}
(a) $\sigma = \sqrt{0.12 \times 0.88/800} = 0.0115$, atau $12.0 \pm 1.1$ poin persentase, dengan
galat relatif $0.0115/0.12 = \qty{9.6}{\percent}$. (b) Titik yang empat kali lebih banyak membuatnya separuh: $0.0057$.
\end{solution}

\begin{solution}{exo:b2:field-lab-methods-2:4}
(a) Kehilangannya \qty{0.0288}{g}, dan $0.0288/0.8000 = \qty{3.60}{\percent}$.
(b) $0.7712/8.7712 = 0.0879$.
(c) Serbuk yang telah dipijarkan: $5.20/0.08792 = \qty{59.1}{\percent}$; serbuk semula:
$59.14 \times (1 - 0.0360) = \qty{57.0}{\percent}$.
\end{solution}

\begin{solution}{exo:b2:field-lab-methods-2:5}
(a) $\sigma = 0.10 \times 0.02 = 0.002$ dan $n = 0.02 \times 0.98/0.002^2 = 4900$ titik.
(b) $\sigma = 0.03$ dan $n = 0.30 \times 0.70/0.03^2 = 233.3$, jadi 234 titik. (c) Untuk galat relatif
yang tetap, $n = (1-p)/(p\,\varepsilon^2)$ tumbuh sebagai $1/p$: mineral aksesori itu menuntut titik 21 kali lebih
banyak, lebih banyak daripada yang biasanya tersedia pada satu sayatan.
\end{solution}

\begin{solution}{exo:b2:field-lab-methods-2:6}
(a) RGM-1: $72.89/73.45 = \qty{99.2}{\percent}$; BCR-2: $54.17/54.1 = \qty{100.1}{\percent}$.
(b) RGM-1: $0.19/\sqrt{20} = 0.042$, dan $-0.56/0.042 = -13.2$; BCR-2: $0.13/\sqrt{20} = 0.029$, dan
$0.07/0.029 = 2.4$. (c) RGM-1 berbias: rata-ratanya meleset dari nilai rekomendasi sebesar 13 kali simpangan
bakunya sendiri, jauh lebih besar daripada yang dimungkinkan oleh kebetulan (jika nilai rekomendasinya benar). BCR-2 berada di
batas yang dapat diberikan oleh kebetulan. RGM-1 tidak kurang presisi: simpangan baku relatifnya,
$0.19/72.89 = \qty{0.26}{\percent}$, dekat dengan simpangan baku relatif BCR-2, \qty{0.24}{\percent}. Masalahnya adalah
kebenaran, bukan presisi.
\end{solution}

\begin{solution}{exo:b2:field-lab-methods-2:7}
(a) $M = 39.098 + 3 \times 24.305 + 26.982 + 3 \times 28.085 + 12 \times 15.999 + 2 \times 1.008 =
\qty{417.25}{g.mol^{-1}}$; kedua OH membentuk satu \ce{H2O} ($18.015$), \qty{4.32}{\percent}.
(b) \ce{K2O} $47.10/417.25 = \qty{11.3}{\percent}$; \ce{MgO} $3 \times 40.304/417.25 =
\qty{29.0}{\percent}$; \ce{Al2O3} $50.98/417.25 = \qty{12.2}{\percent}$; \ce{SiO2} $3 \times
60.083/417.25 = \qty{43.2}{\percent}$; total \qty{95.7}{\percent}. (c) Mikroprobe tidak
mendeteksi hidrogen: analisis mika kehilangan airnya, sekitar \qty{4}{\percent}, dan semestinya berjumlah sekitar
96.
\end{solution}

\begin{solution}{exo:b2:field-lab-methods-2:8}
(a) $\beta = \ln(0.1194/0.11905)/\ln(85.909261/87.905612) = 0.002936/(-0.02297) = -0.1278$.
(b) $0.71062 \times (86.908878/85.909261)^{-0.1278} = 0.70957$. (c) Rasio 86/88 terukur berada di bawah
rasio sebenarnya: isotop berat diunggulkan, dan 87/86 terukur terlalu tinggi.
\end{solution}

\begin{solution}{exo:b2:field-lab-methods-2:9}
(a) $0.1000 \times 15.271 \times 10^{-6}/83.9242 \times 1.3599 \times 10^{-3} = \qty{2.474e-11}{mol}$.
(b) Menurut \eqref{eq:b2:field-lab-methods-2:id}, $N_x = 2.474 \times 10^{-11} \times (732.8 -
0.874)/(0.874 - 0.0568) = \qty{2.216e-8}{mol}$ $^{86}\mathrm{Sr}$; stronsium: $2.216 \times
10^{-8}/0.0986 \times 87.62 = \qty{1.969e-5}{g}$ dalam \qty{0.0500}{g}, \qty{394}{ppm} (393.9).
(c) Perolehan kembali \qty{99.95}{\percent}: metode itu akurat dalam batas pembulatan rasionya.
\end{solution}

\begin{solution}{exo:b2:field-lab-methods-2:10}
(a) $R_m = \sqrt{732.8 \times 0.05680} = 6.45$, tempat $\mathcal M = 1.018$. (b) $\mathcal M = 2.31$
pada 0.1, 1.13 pada 0.5, 1.038 pada 1.664, dan 1.16 pada 100. (c) Di dekat $R_x$ campuran itu hampir tidak berbeda dari
sampel, di dekat $R_s$ dari spike, dan selisih yang membawa informasi hilang dalam
galat $R_m$: $\mathcal M$ divergen di kedua ujung. Di antara keduanya minimumnya sangat datar
($\mathcal M < 1.2$ dari sekitar 0.4 sampai 100), sehingga analis hanya memerlukan orde besaran yang benar
untuk spike, yang dapat diberikan oleh konsentrasi kasar, misalnya dari XRF.
\end{solution}

\begin{solution}{exo:b2:field-lab-methods-2:11}
(a) $\sigma = \sqrt{0.62 \times 0.38/400} = 0.0243$: $62 \pm 2.4$. (b) $(65 - 62)/2.43 = 1.2$
simpangan baku: batuan itu bisa saja granodiorit. (c) $\sigma = 0.015$ memerlukan $0.62 \times
0.38/0.015^2 = 1047.1$, jadi 1048 titik feldspar.
\end{solution}

\begin{solution}{exo:b2:field-lab-methods-2:12}
(a) Setiap titik meleset dengan probabilitas $1 - p$, secara saling bebas: $(1-p)^n$. (b) $0.99^n \le 0.05$
memberikan $n \ge \ln 0.05/\ln 0.99 = 298.1$: 299 titik. (c) $n \ge \ln 0.05/\ln 0.9995 = 5990$: penghitungan
pada kisi adalah alat yang salah; mineral aksesori yang jarang dicari dengan memindai seluruh sayatan.
\end{solution}

\begin{solution}{pb:b2:field-lab-methods-2:1}
\noindent\textbf{Bagian I.}
\textbf{1.} Kuarsa $372/1200 = \qty{31.0}{\percent}$, feldspar alkali 29.0, plagioklas 27.0, biotit
9.0, muskovit 3.0, oksida opak 1.0.
\textbf{2.} $\sqrt{0.31 \times 0.69/1200} = 0.0134$; muskovit: $\sqrt{0.97/(0.03 \times 1200)} =
\qty{16.4}{\percent}$.
\textbf{3.} $M = 9 + 3 + 1 = \qty{13.0}{\percent}$; $M' = 13 - 3 = \qty{10.0}{\percent}$ tanpa
muskovit: leukokratik.
\textbf{4.} $Q + A + P = 1044$ titik: Q $= 35.6$, A $= 33.3$, P $= 31.0$.
\textbf{5.} $324/672 = \qty{48.2}{\percent}$, dengan Q antara 20 dan 60: granit, dan karena antara 35 dan
65, monzogranit.
\textbf{6.} $\sqrt{0.482 \times 0.518/672} = 0.0193$: batas-batas pada 35 dan 65 terletak 6.9 dan 8.7
simpangan baku jauhnya. Nama itu meyakinkan.
\textbf{7.} $0.31 \times 0.69/0.01^2 = 2139$ titik.

\noindent\textbf{Bagian II.}
\textbf{8.} $0.0051/0.8000 = \qty{0.64}{\percent}$.
\textbf{9.} $0.7949/8.7949 = 0.0904$, dan $72 \times 0.0904 = \qty{6.51}{\percent}$ \ce{SiO2} di dalam
kaca: peleburan mengencerkan sampel sebelas kali.
\textbf{10.} $\bar I = 24.91$, $\bar c = 62.283$, $\sum(I - \bar I)^2 = 31.918$, $\sum(I - \bar I)(c -
\bar c) = 80.009$: $a = 2.5067$ per kcps, $b = 62.283 - 2.5067 \times 24.91 = -0.16$.
\textbf{11.} $2.5067 \times 28.63 - 0.16 = \qty{71.61}{\percent}$.
\textbf{12.} BCR-2 $-0.04$, AGV-2 $+0.05$, RGM-1 $-0.01$: semuanya lebih kecil daripada simpangan baku
metode (0.1 sampai 0.2), sehingga garis itu memadai dalam rentang ini.
\textbf{13.} $71.61 + 27.15 + 0.64 = \qty{99.40}{\percent}$, antara 97 dan 101: tidak diulang.
\textbf{14.} $71.61 \times 100/(71.61 + 27.15) = \qty{72.51}{\percent}$.

\noindent\textbf{Bagian III.}
\textbf{15.} $0.0986/0.8258 = 0.1194$: nilai normalisasi itu adalah rasio alami.
\textbf{16.} $\beta = 0.007313/(-0.022972) = -0.3183$.
\textbf{17.} $0.71655 \times (86.908878/85.909261)^{-0.3183} = 0.71392$.
\textbf{18.} Rasio 86/88 terukur menyimpang $0.11853/0.1194 - 1 = \qty{-0.73}{\percent}$ dari nilai sebenarnya
untuk dua satuan massa: sekitar \qty{0.36}{\percent} per satuan massa, menguntungkan isotop-isotop berat.
\textbf{19.} Rasio $^{86}\mathrm{Sr}/^{88}\mathrm{Sr}$ spike adalah 0.7661: rasio 86/88 campuran tidak lagi
0.1194 yang alami, dan menetapkannya pada 0.1194 berarti memperlakukan spike sebagai bias instrumen.

\noindent\textbf{Bagian IV.}
\textbf{20.} Jumlah rasio terhadap 88: $561.4 + 0.7661 + 0.1819 + 1 = 563.348$; $561.4/563.348 =
\qty{99.65}{\percent}$ dan $0.7661/563.348 = \qty{0.13599}{\percent}$.
\textbf{21.} $0.2500 \times 15.271 \times 10^{-6}/83.9242 = \qty{4.549e-8}{mol}$ stronsium, dan
$\times 1.3599 \times 10^{-3} = \qty{6.186e-11}{mol}$ $^{86}\mathrm{Sr}$.
\textbf{22.} $561.4/0.7661 = 732.8$ dan $0.0056/0.0986 = 0.05680$.
\textbf{23.} $N_x = 6.186 \times 10^{-11} \times (732.8 - 1.6812)/(1.6812 - 0.0568) =
\qty{2.784e-8}{mol}$.
\textbf{24.} $\mathcal M = 1.6812 \times 732.74/(731.12 \times 1.6244) = 1.04$; rasio itu menyumbang
$1.04 \times \qty{0.01}{\percent}$, spike $0.031/15.271 = \qty{0.20}{\percent}$; dijumlahkan secara kuadratur
\qty{0.20}{\percent}, didominasi oleh kalibrasi spike.
\textbf{25.} Stronsium: $2.784 \times 10^{-8}/0.0986 = \qty{2.824e-7}{mol}$, yaitu $2.824 \times
10^{-7} \times 87.62 = \qty{2.474e-5}{g}$ dalam \qty{0.1000}{g}. Granit Varden mengandung
\textbf{\qty[text-series-to-math=true]{247}{ppm}} stronsium, $\pm 0.5$.
\end{solution}
